% !TEX root = ../main.tex
\section{Experiment}
\label{sec:experiment}

% New section, 2026-10-11, from the owner's note at the foot of page 12 of the
% printed report ("3 Experiment / 3.1 Scenario"); see handwritten/2026-10-11_main.md.
% Built from the old subsections "Experiment Design" and "Scoring Model
% Calculations" of 03_assumptions_methodology.tex. The scoring calculations were
% AI-drafted on 2026-10-11 and are cut here as the owner marked: equations that
% repeat the OpenAP and Bartel and Young papers are replaced by citations.
% Structure of the experiment: knowledge/models/experiment-definition.md and
% knowledge/models/experiment-options-and-scoring.md; decisions D-013 to D-025.

\plantodo{S11}{\textbf{State the central claim} the experiment explores. Still undecided (see \texttt{experiment-definition.md} \S1 and the hand review notes 1.2, 2.1, 4.1). Answer the objection that the simulation compares two ways of operating and so does not by itself test MBSE.}

\subsection{Scenario}
\label{sec:scenario}

\plantodo{S11}{\textbf{Scenario.} Four aircraft in one invented en-route sector, shown as one snapshot: a crossing pair at the same level and two aircraft in trail 2,000 ft below. State which aircraft belong to the same airline, which is late with connections, and the weight of each. Give the level of the crossing pair and say why it is above FL350. Justify the static snapshot by the controller studies that used one (\cite{kirwan2001coraStrategies}, \cite{fothergill2013resolutionHeuristics}).}

\plantodo{S11 (owner's note, 2026-10-11 print, page 11)}{\textbf{``Justify one scenario (time) and expand as possible.''} Say that one scenario is studied because of the time available in one term, name what a second scenario would add, and list the variations that are run on the one scenario (similar weights, the climb closed, a crossing pair at or below FL350).}

\plantodo{S11 (found while placing the figure)}{\textbf{The drawing and the scenario file disagree on one number.} Figure~\ref{fig:scenario-plan} is labelled ``20 min to WPT''; \texttt{simulation/scenario\_base.yaml} has \texttt{time\_to\_crossing\_min: 12}, meaning the decision is made 12 minutes before the two aircraft reach the crossing. The scoring code does not read that value, so no result changes whichever is kept. Decide which is right and change the other.}

\begin{figure}[htbp]
\centering
\figorbox[0.9\textwidth]{4-ac-scenario_DEN-LAX-crossing.png}{MISSING: \texttt{attachments/4-ac-scenario\_DEN-LAX-crossing.png}}
\caption{The scenario in plan view. AC1 and AC2 cross at the waypoint at FL380; AC3 and AC4 follow AC1's track at FL360.}
\label{fig:scenario-plan}
\end{figure}

\begin{figure}[htbp]
\centering
\figorbox[0.6\textwidth]{4-ac-scenario_altitude_vs_time_plot.png}{MISSING: \texttt{attachments/4-ac-scenario\_altitude\_vs\_time\_plot.png}}
\caption{Level of each aircraft against its time at the waypoint. AC1 and AC2 arrive together at FL380.}
\label{fig:scenario-profile}
\end{figure}

\subsection{Cases and Resolution Options}
\label{sec:experiment-method}

\plantodo{S11}{\textbf{The one thing that differs between cases.} Define cases A to D by what the decider knows. State what ``automated decision'' means for who devises and who issues the clearance (undecided). State that the levers are the same in every case.}

\begin{table}[h]
\centering
\caption{The four cases. Wording to be the author's.}
\label{tab:cases}
\begin{tabular}{@{}llp{6.2cm}l@{}}
\toprule
Case & Decider & Information & Climb rule applies? \\
\midrule
A & Controller, by rule & Type, level, speed, route & Yes \\
B & Automation & Nominal weight per type & Yes \\
C & Automation & Shared weight and cost index & No \\
D & Automation & Shared weight, cost index and delay margin & No \\
\bottomrule
\end{tabular}
\end{table}

\plantodo{S11}{\textbf{Altitude options and the climb rule (assumption adopted from a prior study).} State the three altitude options (2,000 ft above, 2,000 ft below, 4,000 ft below) and the rule that a climb is offered only at or below FL350, citing \cite{engility2014tasarAlaska} (p.~8). Say plainly that this is that study's modeling assumption, made because it did not model aircraft weight; it is not an FAA rule and not observed controller behaviour. Then explain why the rule binds cases A and B and not cases C and D: this is the mechanism by which missing information removes an option.}

\plantodo{S11}{\textbf{Baseline rule for case A.} One descent clearance, citing \cite{rantanen2012conflictManeuvers}. State which aircraft descends and the tie-breaker (undecided), and whether the aircraft is cleared back up.}

\plantodo{S11}{\textbf{Resolution options.} Introduce Table~\ref{tab:options}. For each option give the instructions it costs now, the follow-up it creates, and what the decider must know to use it. Cite \cite{fothergill2013resolutionHeuristics} for the instruction accounting.}

% Moved here from the Results section 2026-10-11: the option set is part of the
% experiment's design, not a result (Claude's suggestion; the owner's print only
% struck "PLACEHOLDER" from the caption).
\begin{table}[h]
\centering
\caption{Resolution options. A dash means no instruction. Full version with instruction counts in Appendix~\ref{app:optimization}.}
\label{tab:options}
\begin{tabular}{@{}lllll@{}}
\toprule
Option & AC1 & AC2 & AC3 & AC4 \\
\midrule
O1 & Descend 2,000 ft & -- & -- & -- \\
O2 & -- & Descend 2,000 ft & -- & -- \\
O3 & Climb 2,000 ft & -- & -- & -- \\
O4 & -- & Climb 2,000 ft & -- & -- \\
O5 & Turn behind AC2 & -- & -- & -- \\
O6 & -- & Turn behind AC1 & -- & -- \\
O7 & Reduce speed & -- & -- & -- \\
O8 & Descend 4,000 ft & -- & -- & -- \\
O9 & Descend 2,000 ft & -- & Slow or turn & --\\
\bottomrule
\end{tabular}
\end{table}

\plantodo{S11}{\textbf{Sources for the option set and the counts} (owner's note, second hand review, 2.4: ``Add these to report!''). Levers: \texttt{experiment-definition.md} \S4, and \cite{engility2014tasarAlaska} (p.~8) for lateral and altitude changes; that study does not use speed. Altitude steps and climb rule: \cite{engility2014tasarAlaska} (p.~8), with the authors' reason. Instruction counts: \cite{fothergill2013resolutionHeuristics} (pp.~72--74). Climbs used less than descents: \cite{rantanen2012conflictManeuvers} (6 of 36). Speed limited to $\pm 0.04$ Mach and built to fail: the author's own figure, to be stated as experience.}

\plantodo{S11}{\textbf{The baseline is not hobbled} (owner's note, second hand review, 3.2). In the later variation at or below FL350 the baseline may climb and does not, because its rule is one descent. Say that this matches the observed preference for descents (\cite{rantanen2012conflictManeuvers}). \textbf{Owner's note, 2026-10-11 print: ``Also note higher alt $\sim$ optimal''} --- say that the higher level is near the optimum for these aircraft (Figure~\ref{fig:fuelflow-altitude}), so a baseline that never climbs is giving something up.}

\subsection{Scoring}
\label{sec:scoring-method}

\plantodo{S11}{\textbf{What the system column is} (owner's notes, second hand review, 5.5 and 6.1). State that it is a measure of the airspace as a whole system of systems, define it by its cost statement (total fuel, total CO2, total airline cost, workload items, largest cost to any one aircraft), and refer back to that statement wherever the column is discussed.}

\plantodo{S11}{\textbf{How each option is scored, one statement per viewpoint.} Aircraft: extra fuel plus cost index times extra time (\cite{mori2022massCruise}, Eq.~1--3). Airline: the sum over its aircraft plus a step cost when a flight misses a connection, plus CO2 converted to dollars. Sector: clearances issued, effort to devise them, and residual risk, reported as three separate items. System: totals stated side by side, with no weighted sum (D-013). Say why no weighted sum: nobody in the model owns the weights.}

\plantodo{S11}{\textbf{Emissions.} State the conversion from fuel to CO2 and its source (\cite{icao2024icecMethodology}). State how CO2 becomes dollars at the airline level and why that route was chosen (credit price, passenger value, or both; undecided), and why the system level keeps CO2 in kg.}

% ---- The text from here to the end of this subsection was AI-drafted on
% 2026-10-11 as a placeholder and pared down the same day as the owner marked
% on the print ("Pare down."). To be checked against
% simulation/experiment_scoring.py and rewritten by the author.
For aircraft $i$ the cost of an option is its extra fuel plus its cost index times its extra time, in kilograms of fuel, the form used in \cite{mori2022massCruise}:
\begin{equation}
c_i = \Delta m_{f,i} + \mathit{CI}_i\, \Delta t_i ,
\label{eq:ac-cost}
\end{equation}
where the extra fuel $\Delta m_{f,i}$ and the extra time $\Delta t_i$ are the option minus the baseline of flying on undisturbed, and $\mathit{CI}_i$ is the cost index of aircraft $i$ in kilograms of fuel per minute.

Each automated case chooses, among the options it is allowed to offer, the one with the lowest estimate
\begin{equation}
J = \sum_i \left( \Delta m_{f,i}(\hat m_i) + \widehat{\mathit{CI}}_i\, \Delta t_i(\hat m_i) \right),
\label{eq:decider}
\end{equation}
where the hats mark what the case believes. Case A does not use Eq.~\ref{eq:decider}; it applies its rule, one descent. Case B uses the type's nominal mass for every aircraft and values time at 1 kg/min. Cases C and D use each aircraft's actual mass and cost index. Case D also rejects an option that delays a flight with connections past its threshold. Whatever a case chooses is then scored at the actual masses, and that score is what the tables report.

Schedule reliability enters as a step: a flight with connections that is delayed past its threshold incurs a cost $K$ on top of Eq.~\ref{eq:ac-cost}. No sourced value of $K$ is held, so the tables report whether the threshold is crossed, not a dollar figure. Carbon dioxide is 3.16 kg per kilogram of fuel \cite{icao2024icecMethodology}. Instructions are counted, now and later, and are not converted to anything.
% ---- end pared-down AI draft

\plantodo{S11 (candidate figure)}{\textbf{Figure~\ref{fig:cost-of-delay} is another paper's figure.} It is Figure 1 of \cite{pilon2019udppMargins}, placed here because it shows the step in the cost of delay that $K$ stands for. Decide: keep it as a reproduction with the credit in the caption, redraw it for this scenario, or remove it.}

\begin{figure}[htbp]
\centering
\figorbox[0.7\textwidth]{single_flight_cost_per_delay.png}{MISSING: \texttt{attachments/single\_flight\_cost\_per\_delay.png}}
\caption{Cost of delay for one flight, with steps where a margin is exceeded. Reproduced from \cite{pilon2019udppMargins}, Figure 1.}
\label{fig:cost-of-delay}
\end{figure}

\subsection{Aircraft Model and Its Simplifications}
\label{sec:scoring-limits}

\plantodo{S11, S12}{\textbf{Truth model (decided, D-017).} OpenAP (open aircraft performance model, TU Delft) scores every option: fuel flow, thrust and drag for the B737-9 at the aircraft's actual mass. State the three selectable fidelities (conceptual: feasibility and instruction counts only; nominal mass; actual mass) and that the decider's case (A to D) and the truth model's fidelity are independent settings. The earlier kinematic model (\texttt{vehicle\_dynamics.py}) is not used by the experiment.}

% ---- AI-drafted 2026-10-11 and cut the same day as the owner marked on the
% print ("Cite another paper for this content. Don't re-create."; "Remove.";
% table "Keep"). The removed equations (true airspeed, drag polar, maximum
% thrust, fuel flow, rate of climb, per-manoeuvre fuel and time) are in git
% history, commit 2b0b038, 03_assumptions_methodology.tex.
Each aircraft is a point mass in steady flight at a constant Mach number. Drag, maximum thrust and fuel flow are taken from OpenAP \cite{sun2020openap}, which uses a two-term drag polar (its Eq.~30 and~31) and, above 30,000 ft, the two-shaft turbofan thrust model of \textcite{bartel2008thrustFuel} (reproduced as its Eq.~22 and~23); those equations are in the OpenAP paper and are not repeated here. The fuel-flow curve is the exception: the installed version of the package (2.6.2) uses a different curve from the one the paper describes, so for fuel flow the paper is the source of the data (the ICAO engine emissions databank) and the package is the source of the curve. Table~\ref{tab:openap-params} lists the values OpenAP holds for the type. Aircraft mass enters through drag: a heavier aircraft needs more lift, and its drag rises with the square of the lift coefficient.

\plantodo{S11, S12 (owner's note, 2026-10-11 print, page 13)}{\textbf{``Cite another paper for this content. Don't re-create.''} The OpenAP paper is now held (\texttt{evidence/sources/sun2020openap\_aerospace-07-00104.pdf}) and was read in full by Claude on 2026-10-11; the equation numbers above are from that reading. The author has not read it. Bartel and Young is not held, and is cited here only through the OpenAP paper, which reproduces its equations. Read Sections 5 and 7 of the OpenAP paper (pp.~6 to 17) before submission. What the paper says about its own accuracy is in \texttt{evidence/literature-notes/summaries/4 - sun2020openap.md} and belongs in the discussion of the truth model's uncertainty (Section~\ref{sec:experiment-results}).}

\begin{table}[h]
\centering
\small
\caption{Values OpenAP 2.6.2 uses for the B737-9 (type code B39M), read from the installed package on 2026-10-11.}
\label{tab:openap-params}
\begin{tabular}{@{}p{5.6cm}p{3.6cm}p{5.6cm}@{}}
\toprule
Quantity & Value & Note \\
\midrule
Wing area $S$ & 124.6 m$^2$ & type data \\
Zero-lift drag coefficient $C_{D0}$ & 0.020 & 737 MAX 8 polar, borrowed \\
Induced drag factor $k$ & 0.042 & 737 MAX 8 polar, borrowed \\
Engines $n$, type & 2, LEAP-1B & type data \\
Static thrust per engine $T_0$ & 111,300 N & engine data \\
Reference thrust $T_{cr}$ & 46,300 N (both engines) & $0.2\,T_0 + 890$ N per engine \\
Reference cruise $M_{cr}$, altitude & 0.79, 11,000 m & type data \\
Take-off fuel flow $s$ & 0.877 kg/s per engine & ICAO databank \\
Fuel curve $c_1$, $c_2$, $c_3$ & 0.9376, 1.9768, 1.3955 & package default curve, not fitted to this type \\
\bottomrule
\end{tabular}
\end{table}

A descent is flown at 1,500 ft/min. A climb is flown at maximum thrust, at the rate that thrust allows and no more than 1,500 ft/min (D-024). A level is available to an aircraft only if it would still have 150 ft/min of climb left there. Cases A and B do not know the mass and so cannot make that test; they offer a climb only at or below FL350. Cases C and D make it at the aircraft's actual mass.

% Figure moved here from the Results section and its caption cut, per the
% owner's notes on the 2026-10-11 print, page 20 ("Move figure to an earlier
% section."; title: "Move to caption"; bracketed caption text: "Move to
% discussion."). The sentence after the figure is the text moved out of the caption.
\begin{figure}[htbp]
\centering
\figorbox[\textwidth]{b739_envelope.png}{PLACEHOLDER: run \texttt{simulation/b739\_envelope.py} to produce \texttt{report/figures/b739\_envelope.png}}
\caption{The heavy B737-9 (71.0 t) has little room left to change speed or level at FL380. Left: cruise envelope, with three lines of constant rate of climb left: 0 ft/min (thrust equals drag), 100 ft/min (the service ceiling, \cite{faa2015ac61107b}) and 150 ft/min (the operator margin the scenario uses). Right: width of the usable speed band and rate of climb left at Mach 0.80, by flight level.}
\label{fig:envelope}
\end{figure}

Figure~\ref{fig:envelope} is drawn assuming lift equals weight. Its low-speed boundary is 1.3 times the stall speed at an assumed clean maximum lift coefficient, and its high-speed boundary is the maximum operating Mach number. The aircraft's present level and Mach number are marked, and each altitude and speed instruction the scenario offers it is marked as inside or outside the region with 150 ft/min of climb left.

The model makes the following simplifications.
\begin{enumerate}
\item \textbf{Standard atmosphere, no wind, lift equal to weight.} Time differences come only from true airspeed and added track distance.
\item \textbf{No compressibility drag, and the climb left is sensitive to it.} The drag used here has no Mach-number term. The OpenAP paper does include one, a wave-drag term above a critical Mach number, and says the effect ``should not be ignored'' there \cite{sun2020openap} (Section~5.5); the installed package leaves the term off unless asked and marks it experimental, and the experiment runs with it off (D-024). The FAA classes flight from Mach 0.75 to 1.20 as transonic \cite{faa2015ac61107b} (para.~1-4e), so the scenario's Mach 0.80 is at the low end of that range. With OpenAP's experimental wave-drag term switched on, drag at Mach 0.80 is 2 to 4 percent higher. That is 1.3 kN for the heavy aircraft at FL380, against 2.7 kN of excess thrust, so its climb left falls from 181 to 93 ft/min, under the 100 ft/min of the service ceiling. The light aircraft at FL400 falls from 322 to 239 ft/min and still clears 150. The weights in the scenario were set with the term off. Not decided: whether to switch the term on and set the weights again.
\item \textbf{The drag polar is the 737 MAX 8's.} OpenAP has none for the 737 MAX 9 (D-017).
\item \textbf{The fuel curve is the package default.} OpenAP has no fuel model fitted to the 737 MAX 9 or MAX 8; it uses its default coefficients scaled by the LEAP-1B take-off fuel flow.
\item \textbf{The climb margins.} The service ceiling is the level with 100 ft/min of climb left \cite{faa2015ac61107b}. The 150 ft/min the experiment requires is the margin the author has seen airlines keep for clear air turbulence and other unknowns (industry experience, no document held; D-023).
\item \textbf{One climb rate per climb.} The rate is taken at the mid-altitude and held; the climb is not integrated. A climb is not checked against the time left before the crossing; the longest in the base scenario is the heavy aircraft's 4,000 ft climb back after option O8, about 12 min, which happens after the conflict.
\item \textbf{Mass is constant.} Fuel burned during the manoeuvre does not reduce the mass.
\item \textbf{The low-speed limit rests on an assumed maximum lift coefficient} ($C_{L\max} = 1.45$) and is a stall margin, not a buffet boundary.
\item \textbf{A speed change on the crossing pair is not calculated.} It is marked unable to separate the pair by construction (D-014).
\end{enumerate}
% ---- end AI draft

\plantodo{S11}{\textbf{Abstractions made and why.} One short table: each simplification, the reason, and the risk. A draft list is in \texttt{experiment-options-and-scoring.md} \S6.}

\subsection{Worked Example: Option O4, the Light Aircraft Climbs 2,000 ft}
\label{sec:worked-example}

% ---- AI-drafted 2026-10-11; kept and rewritten the same day to stand without
% the removed equations (owner's note: "Keep. Update w/ eqn removals &
% simplification of explanations").
The light aircraft, AC2, is at 62,980 kg, FL380 and Mach 0.80. At FL400 it would have 322 ft/min of climb left, more than the 150 ft/min required, so the level is available to a case that knows the mass.

At maximum thrust it climbs at 409 ft/min, so the 2,000 ft take 294 seconds. OpenAP gives a fuel flow of 0.5572 kg/s at FL380, 0.6321 kg/s in the climb, and 0.5465 kg/s at FL400. The aircraft then holds FL400 for 15 minutes (900 seconds). The extra fuel is what the climb costs less what the higher level saves:
\begin{equation*}
\Delta m_f = (0.6321 - 0.5572)(294) + (0.5465 - 0.5572)(900) = 22.0 - 9.6 = 12.4 \text{ kg}.
\end{equation*}
The extra time is zero: at these levels the standard temperature is constant, so Mach 0.80 is the same true airspeed at FL380 and FL400. With no extra time, the cost index does not enter, and the cost of the option to AC2 is 12.4 kg. This is the 12 kg, rounded, in the case C row of Table~\ref{tab:outcomes}.
% ---- end AI draft
